\subsection*{\texttt{span} and \texttt{span(i)}}

\texttt{tensor.span} gives the full shape of a tensor. \texttt{tensor.span(i)} gives one dimension.

\begin{lstlisting}[language=croqtile,caption={span and span(i)}]
s32 [lhs.span(0), rhs.span(1)] output;
\end{lstlisting}

This comes up constantly in data movement code because output shapes often come from different inputs, and moved tiles frequently need to preserve or recompute shape information.

\begin{figure}[htbp]\centering
\includegraphics[width=0.86\linewidth,height=0.42\textheight,keepaspectratio]{f04-fig_span_light.png}
\end{figure}
\begin{center}\small\emph{Use \texttt{span} for the whole shape, \texttt{span(i)} for a single axis.}\end{center}
